\documentclass[11pt,a4paper,oneside]{book}
\usepackage{mathscommun}
\setniveau{Terminale}
\renewcommand{\ReferenceProgramme}{Spécialité (2026--2027 et 2027--2028) et mathématiques expertes}
\renewcommand{\EditionDocument}{Édition du 24 octobre 2026}
\hypersetup{pdftitle={Terminale - Exercices et problèmes}}
\begin{document}
\titrecollection{Exercices et problèmes}{Spécialité, mathématiques expertes et fondements pour le supérieur.}{Un parcours démontré : maîtriser les définitions, choisir les méthodes, contrôler les hypothèses et préparer la licence avant les spécialisations de master.}
\frontmatter
\chapter*{Mode d'emploi}
Ce volume suit le parcours de terminale générale avec spécialité mathématiques et option mathématiques expertes. Il ne constitue pas un cours des voies technologiques ni de l'option mathématiques complémentaires. La table de couverture du guide distingue les textes applicables en 2026--2027 et à partir de 2027--2028.

Les chapitres 1 à 24 forment le parcours de spécialité, avec les sections signalées en approfondissement, notamment les sphères et les barycentres du chapitre 8. Les chapitres 25 à 35 couvrent l'option expertes. Le chapitre 36 est une passerelle : il donne les fondements de plusieurs preuves et ne prétend pas remplacer une licence.

Les rappels de première sont repris au point d'usage. Pour une remise à niveau plus longue, revenir aux chapitres correspondants du volume de première : raisonnement (1), suites (4--7), dérivation (8--11), exponentielle (12), trigonométrie (14--15), produit scalaire (16), probabilités (19--23).

Lire une première fois pour comprendre les objets ; refaire ensuite les preuves sans support ; enfin chercher les exercices et les synthèses avant de consulter les corrigés. Chaque chapitre possède six exercices numérotés de la même façon dans les deux recueils. Les indications de difficulté distinguent les tâches de base, de raisonnement et d'approfondissement.

Les fondements admis sont recensés dans le guide. Les quantificateurs, la complétude et les notations abstraites supplémentaires sont explicités pour préparer le supérieur ; leur maîtrise complète dépasse les exigences scolaires.
\tableofcontents
\mainmatter
\chapter{Langage, fonctions et calcul : les bases}
\section*{Exercice 1.1 -- Nier une affirmation}
[Base] Nier : pour tout réel $x$, il existe un entier $n$ tel que $n>x$. L'affirmation initiale est-elle vraie ?

\section*{Exercice 1.2 -- Domaine avant calcul}
[Base] Déterminer le domaine de $f(x)=\ln(x-1)/\sqrt{4-x}$ puis résoudre $f(x)=0$.

\section*{Exercice 1.3 -- Une équivalence à corriger}
[Raisonnement] Résoudre $\sqrt{x+6}=x$. Expliquer pourquoi une élévation au carré seule ne suffit pas.

\section*{Exercice 1.4 -- Composer}
[Base] Pour $f(x)=x^2$ et $g(x)=x+1$, calculer $g\circ f$ et $f\circ g$. Sont-elles égales ?

\section*{Exercice 1.5 -- Bijection et restriction}
[Raisonnement] Donner un domaine et un ensemble d'arrivée pour rendre $x\mapsto x^2$ bijective, puis écrire sa réciproque.

\section*{Exercice 1.6 -- Égalité d’ensembles}
[Preuve] Pour $f:E\to F$ et $A,B\subset F$, montrer $f^{-1}(A\cup B)=f^{-1}(A)\cup f^{-1}(B)$. Ici $f^{-1}(A)$ désigne une image réciproque, même sans bijection.

\chapter{Récurrence et sommes}
\section*{Exercice 2.1 -- Somme impaire}
[Preuve] Démontrer $1+3+\cdots+(2n-1)=n^2$ pour $n\ge1$.

\section*{Exercice 2.2 -- Somme géométrique}
[Base] Calculer $\sum_{k=2}^6 3^k$.

\section*{Exercice 2.3 -- Télescopage}
[Méthode] Calculer $\sum_{k=1}^n1/((k+1)(k+2))$.

\section*{Exercice 2.4 -- Une récurrence fausse}
[Critique] Quel défaut contient : « si $u_n=1$, alors $u_{n+1}=1$, donc tous les termes valent 1 », pour $u_0=2$, $u_{n+1}=u_n^2$ ?

\section*{Exercice 2.5 -- Bernoulli appliqué}
[Base] Donner une borne inférieure simple de $1{,}01^{100}$ sans calculatrice.

\section*{Exercice 2.6 -- Divisibilité par récurrence}
[Preuve] Montrer $3\mid(4^n-1)$ pour tout $n\ge0$.

\chapter{Dénombrer sans compter deux fois}
\section*{Exercice 3.1 -- Codes}
[Base] Compter les codes de quatre chiffres distincts, zéros initiaux autorisés. Puis les nombres de quatre chiffres distincts.

\section*{Exercice 3.2 -- Comité}
[Base] Parmi 10 personnes, former un comité de 4 dont deux personnes fixées doivent faire partie.

\section*{Exercice 3.3 -- Au moins un}
[Méthode] Parmi 5 femmes et 4 hommes, combien de groupes de 3 comportent au moins une femme ?

\section*{Exercice 3.4 -- Réunion}
[Base] Sur 40 élèves, 22 font A, 18 font B et 9 font les deux. Compter ceux qui font exactement une activité, puis aucune.

\section*{Exercice 3.5 -- Étoiles et séparateurs}
[Approfondissement] Compter les triplets d'entiers positifs ou nuls de somme 6.

\section*{Exercice 3.6 -- Sous-ensembles}
[Preuve] Dans un ensemble de $n\ge1$ éléments, combien de parties contiennent un élément fixé ?

\chapter{Coefficients binomiaux et formule du binôme}
\section*{Exercice 4.1 -- Pascal}
[Base] Construire la ligne $n=5$ du triangle de Pascal, puis vérifier sa somme.

\section*{Exercice 4.2 -- Coefficient}
[Base] Trouver le coefficient de $x^2$ dans $(1-2x)^5$.

\section*{Exercice 4.3 -- Somme alternée}
[Preuve] Pour $n\ge1$, interpréter combinatoirement $\sum_{k=0}^n(-1)^k\binom nk=0$.

\section*{Exercice 4.4 -- Élément distingué}
[Preuve] Démontrer $\sum_{k=0}^n k\binom nk=n2^{n-1}$ pour $n\ge1$.

\section*{Exercice 4.5 -- Somme de carrés}
[Approfondissement] Montrer $\sum_{k=0}^n\binom nk^2=\binom{2n}n$.

\section*{Exercice 4.6 -- Binôme et matrices}
[Critique] Développer $(A+B)^2$ pour deux matrices carrées et donner la condition qui permet la formule scalaire.

\chapter{Vecteurs, bases et repères de l’espace}
\section*{Exercice 5.1 -- Cube}
[Base] Dans le cube du cours, calculer $\overrightarrow{EC}$ et les coordonnées du milieu de $[BH]$.

\section*{Exercice 5.2 -- Combinaison}
[Base] Écrire $(3,1,2)$ dans la base $u=(1,0,0)$, $v=(1,1,0)$, $w=(0,0,1)$.

\section*{Exercice 5.3 -- Plan ou droite}
[Critique] L'ensemble $A+su+tv$ est-il toujours un plan ? Donner un contre-exemple.

\section*{Exercice 5.4 -- Coplanarité}
[Méthode] Les vecteurs $(1,0,1)$, $(0,1,1)$ et $(2,3,5)$ sont-ils coplanaires ?

\section*{Exercice 5.5 -- Droites gauches}
[Preuve] Justifier que $(t,0,0)$ et $(0,s,1)$ sont non coplanaires.

\section*{Exercice 5.6 -- Unicité}
[Preuve] Montrer que trois vecteurs formant une base ne peuvent satisfaire $u+2v-3w=0$.

\chapter{Produit scalaire, orthogonalité et projections}
\section*{Exercice 6.1 -- Angle}
[Base] Calculer l'angle entre $(1,1,0)$ et $(0,1,1)$.

\section*{Exercice 6.2 -- Projection sur une droite}
[Méthode] Projeter $M=(2,0,1)$ sur la droite passant par l'origine de directeur $(1,1,0)$.

\section*{Exercice 6.3 -- Normalité}
[Preuve] Montrer que $(1,1,1)$ est normal au plan engendré par $(1,-1,0)$ et $(0,1,-1)$.

\section*{Exercice 6.4 -- Deux plans}
[Base] Les plans de normales $(1,2,0)$ et $(2,-1,3)$ sont-ils perpendiculaires ?

\section*{Exercice 6.5 -- Minimum}
[Preuve] Pour un projeté $H$ de $M$ sur un plan, prouver que $K\ne H$ dans le plan donne $MK>MH$.

\section*{Exercice 6.6 -- Inégalité triangulaire}
[Approfondissement] Déduire $\|u+v\|\le\|u\|+\|v\|$ de Cauchy--Schwarz.

\chapter{Droites, plans et systèmes linéaires}
\section*{Exercice 7.1 -- Équation de plan}
[Base] Donner le plan passant par $(1,0,2)$ et normal à $(2,-1,3)$.

\section*{Exercice 7.2 -- Intersection}
[Méthode] Intersecter $(x,y,z)=(t,1+t,2-t)$ avec $x+y+z=6$.

\section*{Exercice 7.3 -- Projection sur un plan}
[Méthode] Projeter $(1,2,3)$ sur $x+y+z=0$ et calculer la distance.

\section*{Exercice 7.4 -- Intersection de plans}
[Méthode] Décrire $x+y+z=1$ et $x-y+z=3$.

\section*{Exercice 7.5 -- Paramètres indépendants}
[Critique] Les droites $(t,0,0)$ et $(1+s,s,0)$ se rencontrent-elles ?

\section*{Exercice 7.6 -- Dégénérescence}
[Base] Interpréter $0x+0y+0z=0$ puis $0x+0y+0z=1$.

\chapter{Sphères, barycentres et mesures}
\section*{Exercice 8.1 -- Sphère et plan}
[Base] Intersecter la sphère de centre $O$, rayon 5, avec le plan $z=3$. Déterminer aussi le plan des points équidistants de $A=(0,0,0)$ et $B=(2,0,0)$.

\section*{Exercice 8.2 -- Tangence}
[Base] Donner le plan tangent à la sphère unité en $(0,0,1)$.

\section*{Exercice 8.3 -- Barycentre}
[Base] Calculer le barycentre de $A=(0,0,0)$ et $B=(3,0,0)$ affectés des poids $2$ et $1$.

\section*{Exercice 8.4 -- Leibniz}
[Preuve] Si $I$ est milieu de $[AB]$, démontrer $MA^2+MB^2=2MI^2+AB^2/2$.

\section*{Exercice 8.5 -- Volume}
[Méthode] Calculer le volume du tétraèdre $O,(2,0,0),(0,3,0),(0,0,4)$.

\section*{Exercice 8.6 -- Sphère circonscrite}
[Approfondissement] Trouver le centre de la sphère passant par $O,(2,0,0),(0,2,0),(0,0,2)$.

\chapter{Comprendre la limite d’une suite}
\section*{Exercice 9.1 -- Définition et rang}
[Base] Pour $u_n=3/(n+1)$, trouver un rang à partir duquel $|u_n|<0{,}01$.

\section*{Exercice 9.2 -- Bornée mais divergente}
[Raisonnement] Montrer que $u_n=(-1)^n+1/n$, $n\ge1$, ne converge pas.

\section*{Exercice 9.3 -- Gendarmes}
[Méthode] Étudier $\sin n/n$, $n\ge1$.

\section*{Exercice 9.4 -- Non majorée}
[Critique] Donner une suite positive non majorée ne tendant pas vers $+\infty$.

\section*{Exercice 9.5 -- Ordre et limite}
[Preuve] Si $u_n\to\ell>0$, montrer que $u_n>0$ à partir d'un rang.

\section*{Exercice 9.6 -- Négation de convergence}
[Approfondissement] Écrire la négation exacte de $u_n\to0$.

\chapter{Opérations sur les limites et suites géométriques}
\section*{Exercice 10.1 -- Terme dominant}
[Base] Calculer la limite de $(5n^2+2n)/(n^2-3)$ à l'infini.

\section*{Exercice 10.2 -- Conjugué}
[Méthode] Calculer $\lim(\sqrt{n^2+4n}-n)$.

\section*{Exercice 10.3 -- Géométriques}
[Base] Étudier $(0{,}7)^n$, $(-0{,}7)^n$, $(-2)^n$.

\section*{Exercice 10.4 -- Suite affine}
[Méthode] Résoudre $u_0=1$, $u_{n+1}=u_n/2+3$.

\section*{Exercice 10.5 -- Indétermination}
[Critique] Construire deux quotients de suites tendant toutes vers zéro, l'un tendant vers 2 et l'autre vers $+\infty$.

\section*{Exercice 10.6 -- Somme infinie}
[Approfondissement] Définir puis calculer $\sum_{k=1}^{\infty}3^{-k}$.

\chapter{Suites monotones, points fixes et suites adjacentes}
\section*{Exercice 11.1 -- Suite racine}
[Méthode] Étudier $u_0=1$, $u_{n+1}=\sqrt{3+u_n}$.

\section*{Exercice 11.2 -- Erreur contractante}
[Base] Si $|u_n-L|\le8(1/2)^n$, donner un rang garantissant une erreur au plus $1/100$.

\section*{Exercice 11.3 -- Point fixe trompeur}
[Critique] Pour $u_0=1$ et $u_{n+1}=-u_n$, résoudre l'équation du point fixe et comparer au comportement réel.

\section*{Exercice 11.4 -- Adjacentes}
[Preuve] Pour $a_n=1-1/(n+1)$, $b_n=1+1/(n+1)$, vérifier l'adjacence et encadrer leur limite.

\section*{Exercice 11.5 -- Héron pour 3}
[Méthode] Partant de $u_0=2$, calculer deux itérations de Héron pour $\sqrt3$ et expliquer la convergence.

\section*{Exercice 11.6 -- Limite ou majorant}
[Critique] Une suite croissante est majorée par 100. Peut-on conclure qu'elle tend vers 100 ?

\chapter{Limites de fonctions et croissances comparées}
\section*{Exercice 12.1 -- Limites latérales}
[Base] Étudier $1/(x-2)$ près de 2.

\section*{Exercice 12.2 -- Signe à l’infini}
[Méthode] Calculer les limites de $\sqrt{x^2+4}/x$ aux deux infinis.

\section*{Exercice 12.3 -- Exponentielle dominante}
[Base] Calculer $\lim_{x\to+\infty}x^5e^{-2x}$.

\section*{Exercice 12.4 -- Asymptote oblique}
[Méthode] Étudier les asymptotes de $(x^2+2)/(x-1)$.

\section*{Exercice 12.5 -- Limite et valeur}
[Critique] Poser $f(x)=(x^2-1)/(x-1)$ pour $x\ne1$ et $f(1)=0$. Quelle est sa limite en 1 ? Est-elle continue ?

\section*{Exercice 12.6 -- Une oscillation écrasée}
[Preuve] Montrer que $x\sin(1/x)\to0$ quand $x\to0$.

\chapter{Continuité, valeurs intermédiaires et équations}
\section*{Exercice 13.1 -- Existence et unicité}
[Méthode] Montrer que $x^3+2x-2=0$ a une unique solution réelle, située dans $]0,1[$.

\section*{Exercice 13.2 -- Dichotomie}
[Base] Combien de dichotomies sur $[0,1]$ garantissent que le milieu approche une racine à moins de $10^{-3}$ ?

\section*{Exercice 13.3 -- Prolongement}
[Méthode] Prolonger continûment $(e^x-1)/x$ en 0.

\section*{Exercice 13.4 -- TVI sans monotonie}
[Critique] L'équation $x^2=1/4$ sur $[-1,1]$ a-t-elle une solution unique ?

\section*{Exercice 13.5 -- Image d’une suite}
[Preuve] Si $u_n\to2$, calculer la limite de $\sqrt{1+u_n^2}$.

\section*{Exercice 13.6 -- Additivité continue}
[Approfondissement] Si $f$ est continue et additive avec $f(1)=3$, déterminer $f(\sqrt2)$.

\chapter{Dériver, composer et contrôler les variations}
\section*{Exercice 14.1 -- Chaîne}
[Base] Dériver $e^{3x^2-1}$ et $\ln(1+x^2)$.

\section*{Exercice 14.2 -- Quotient}
[Base] Étudier les variations de $x/(1+x^2)$.

\section*{Exercice 14.3 -- Tangente}
[Méthode] Donner la tangente à $\ln x$ en 1 et expliquer son rôle d'approximation.

\section*{Exercice 14.4 -- Dérivée nulle}
[Critique] Donner une fonction dérivable avec $f'(0)=0$ sans extremum en 0.

\section*{Exercice 14.5 -- Borne d’erreur}
[Preuve] Montrer $|\sin x-\sin y|\le|x-y|$.

\section*{Exercice 14.6 -- Optimisation}
[Synthèse courte] Un rectangle a aire 36. Montrer que son périmètre minimal vaut 24.

\chapter{Convexité, tangentes et inégalités}
\section*{Exercice 15.1 -- Convexe mais décroissante}
[Base] Étudier variations et convexité de $e^{-x}$.

\section*{Exercice 15.2 -- Inflexion}
[Méthode] Déterminer les points d'inflexion de $f(x)=x^4-2x^2$.

\section*{Exercice 15.3 -- Tangente logarithmique}
[Preuve] Démontrer $\ln x\le x-1$ pour $x>0$, avec le cas d'égalité.

\section*{Exercice 15.4 -- Milieu}
[Preuve] Montrer $e^{(a+b)/2}\le(e^a+e^b)/2$.

\section*{Exercice 15.5 -- Un faux critère}
[Critique] Expliquer pourquoi $f''(0)=0$ ne suffit pas à produire une inflexion.

\section*{Exercice 15.6 -- Lorenz}
[Application] Pour les revenus fictifs $1,1,2,6$, construire les points d'une courbe de Lorenz et vérifier la convexité.

\chapter{Logarithme, puissances et modèles multiplicatifs}
\section*{Exercice 16.1 -- Domaine et produit}
[Base] Résoudre $\ln(x-2)+\ln x=\ln3$.

\section*{Exercice 16.2 -- Seuil}
[Méthode] Trouver le plus petit entier $n$ avec $0{,}8^n<0{,}1$.

\section*{Exercice 16.3 -- Limite en zéro}
[Base] Calculer $\lim_{x\to0^+}x^2\ln x$.

\section*{Exercice 16.4 -- Puissance réelle}
[Approfondissement] Dériver $x^{\sqrt2}$ sur $]0,+\infty[$ en revenant à sa définition.

\section*{Exercice 16.5 -- Logarithme d’une somme}
[Critique] Résoudre $\ln(x+1)=\ln x+\ln2$ sans utiliser de règle fausse.

\section*{Exercice 16.6 -- Une limite exponentielle}
[Preuve] Montrer $(1+2/n)^n\to e^2$.

\chapter{Sinus, cosinus et oscillations}
\section*{Exercice 17.1 -- Angles et signes}
[Base] Calculer $\cos(5\pi/3)$, $\sin(-\pi/6)$ et résoudre $\cos x=1/2$ dans $[-\pi,\pi]$.

\section*{Exercice 17.2 -- Inéquation}
[Méthode] Résoudre $\cos x\ge-1/2$ dans $[-\pi,\pi]$.

\section*{Exercice 17.3 -- Composition}
[Base] Dériver $\sin(3x^2)$ et étudier $1+\cos(2x)$ sur $[0,\pi]$.

\section*{Exercice 17.4 -- Rectangle inscrit}
[Modélisation] Un rectangle est inscrit dans un cercle de rayon $R>0$. En paramétrant ses demi-côtés par $R\cos t,R\sin t$, $0<t<\pi/2$, maximiser son aire.

\section*{Exercice 17.5 -- Limite trigonométrique}
[Preuve] Montrer $\lim_{x\to0}\sin(5x)/x=5$ et $\lim_{x\to0}(1-\cos x)/x^2=1/2$.

\section*{Exercice 17.6 -- Tangente et domaine}
[Approfondissement] Définir $\tan x$, dériver cette fonction et expliquer pourquoi on ne peut pas conclure qu'elle est croissante sur tout son domaine.

\chapter{Primitives et équations différentielles}
\section*{Exercice 18.1 -- Primitives et domaines}
[Base] Donner les primitives de $3x^2-2/x$ sur $]0,+\infty[$, puis sur $]-\infty,0[$.

\section*{Exercice 18.2 -- Condition initiale}
[Méthode] Résoudre $y'=2y-6$, $y(0)=5$.

\section*{Exercice 18.3 -- Second membre}
[Méthode] Résoudre $y'=y+x$ en cherchant une solution particulière affine.

\section*{Exercice 18.4 -- Refroidissement fictif}
[Modélisation] Dans un modèle $T'=-0{,}2(T-20)$, $t$ est en minutes et $T(0)=80$ en degrés Celsius. Quand le modèle prévoit-il $T=30$ ?

\section*{Exercice 18.5 -- Euler instable}
[Critique] Pour $y'=-4y$, $y(0)=1$, comparer Euler de pas $h=1$ à la solution exacte.

\section*{Exercice 18.6 -- Logistique}
[Approfondissement] Vérifier que $y(t)=K/(1+Ae^{-rt})$ satisfait $y'=ry(1-y/K)$ pour $K,r>0$ et $A>0$. Donner sa limite et interpréter $K$.

\chapter{Intégrale et théorème fondamental de l’analyse}
\section*{Exercice 19.1 -- Aire signée}
[Base] Calculer $\int_{-1}^1x\,dx$ puis l'aire entre sa courbe et l'axe sur ce segment.

\section*{Exercice 19.2 -- Borne variable}
[Méthode] Dériver $F(x)=\int_1^{x^2}e^{-t^2}\,dt$.

\section*{Exercice 19.3 -- Valeur moyenne}
[Application] Calculer la valeur moyenne de $t^2$ sur $[0,3]$ et trouver un point où elle est atteinte.

\section*{Exercice 19.4 -- Positivité stricte}
[Preuve] Si $f$ est continue, positive ou nulle sur $[a,b]$, $a<b$, et $\int_a^bf=0$, montrer que $f$ est identiquement nulle.

\section*{Exercice 19.5 -- Orientation}
[Base] Sachant $\int_0^2f=5$ et $\int_0^5f=1$, calculer $\int_2^5f$ et $\int_5^2f$.

\section*{Exercice 19.6 -- Comparaison}
[Raisonnement] Pour $0\le x\le1$, encadrer $\int_0^1e^{-x^2}\,dx$ avec des constantes simples, puis améliorer la borne inférieure par $e^{-u}\ge1-u$.

\chapter{Calcul intégral, encadrements et suites d’intégrales}
\section*{Exercice 20.1 -- Intégration par parties}
[Méthode] Calculer $\int_0^1xe^x\,dx$.

\section*{Exercice 20.2 -- Aire entre deux courbes}
[Méthode] Calculer l'aire comprise entre $y=x$ et $y=x^2$ de 0 à 1.

\section*{Exercice 20.3 -- Suite d’intégrales}
[Preuve] Pour $I_n=\int_0^1x^ne^{-x}\,dx$, prouver que $I_n\to0$ et établir $I_n=nI_{n-1}-e^{-1}$ pour $n\ge1$.

\section*{Exercice 20.4 -- Somme de Riemann}
[Approfondissement] Calculer $\lim_{n\to\infty}\frac1n\sum_{k=1}^n\frac1{1+k/n}$.

\section*{Exercice 20.5 -- Trapèzes et convexité}
[Raisonnement] Pour une fonction convexe sur $[a,b]$, montrer $\int_a^bf\le(b-a)(f(a)+f(b))/2$.

\section*{Exercice 20.6 -- Série harmonique}
[Preuve] Utiliser une intégrale pour montrer que $H_n=\sum_{k=1}^n1/k$ diverge.

\chapter{Expériences aléatoires et loi binomiale}
\section*{Exercice 21.1 -- Reconnaître le modèle}
[Critique] On tire 5 cartes sans remise d'un jeu de 32 cartes. Le nombre d'as suit-il une loi binomiale ?

\section*{Exercice 21.2 -- Probabilité exacte}
[Base] Si $X\sim\mathcal B(5,1/3)$, calculer $P(X=2)$ et $P(X\ge1)$.

\section*{Exercice 21.3 -- Atteindre une fiabilité}
[Méthode] Des essais indépendants réussissent avec probabilité $0{,}2$. Combien en faut-il pour avoir au moins une réussite avec probabilité au moins $0{,}95$ ?

\section*{Exercice 21.4 -- Intervalle et quantiles discrets}
[Méthode] Pour $X\sim\mathcal B(10,1/2)$, calculer $P(2\le X\le8)$. Trouver le plus petit $b$ tel que $P(X\le b)\ge0{,}95$.

\section*{Exercice 21.5 -- Indépendance par paires}
[Preuve] Deux bits indépendants équiprobables sont tirés. Montrer que les événements « premier bit égal à 1 », « second égal à 1 » et « bits égaux » sont indépendants deux à deux, mais pas mutuellement.

\section*{Exercice 21.6 -- Optimiser un paramètre}
[Raisonnement] Pour $X\sim\mathcal B(n,p)$ et $0<k<n$ fixé, déterminer $p\in[0,1]$ maximisant $P(X=k)$.

\chapter{Espérance, variance et sommes de variables}
\section*{Exercice 22.1 -- Calcul direct}
[Base] $X$ prend les valeurs $-1,0,2$ avec probabilités $1/4,1/2,1/4$. Calculer son espérance et sa variance.

\section*{Exercice 22.2 -- Transformation affine}
[Base] Sachant $E(X)=3$, $V(X)=4$, calculer l'espérance et l'écart type de $Y=5-2X$.

\section*{Exercice 22.3 -- Somme dépendante}
[Critique] Soit $X$ une Bernoulli de paramètre $1/2$ et $Y=1-X$. Comparer $V(X+Y)$ et $V(X)+V(Y)$.

\section*{Exercice 22.4 -- Binomiale}
[Méthode] Pour $X\sim\mathcal B(100,0{,}3)$, calculer $E(X)$, $V(X)$ et la variance de $X/100$.

\section*{Exercice 22.5 -- Sans indépendance}
[Preuve] Dans une classe fictive de 20 élèves, chacun a probabilité $1/5$ d'oublier un livre. Trouver l'espérance du nombre d'oublis sans supposer l'indépendance.

\section*{Exercice 22.6 -- Moindres carrés}
[Approfondissement] Montrer que $E((X-a)^2)=V(X)+(E(X)-a)^2$. Quel réel $a$ minimise cette quantité ?

\chapter{Concentration et loi des grands nombres}
\section*{Exercice 23.1 -- Bienaymé--Tchebychev}
[Base] Une variable a moyenne 50 et variance 9. Majorer $P(|X-50|\ge6)$.

\section*{Exercice 23.2 -- Taille d’échantillon}
[Méthode] Pour une fréquence de Bernoulli indépendantes de paramètre inconnu, garantir $P(|F_n-p|\ge0{,}05)\le0{,}01$ par Tchebychev.

\section*{Exercice 23.3 -- Rôle de l’indépendance}
[Critique] Si $X_1=\cdots=X_n=X$ avec $X$ Bernoulli de paramètre $1/2$, la moyenne se concentre-t-elle autour de $1/2$ ?

\section*{Exercice 23.4 -- Convergence en probabilité}
[Preuve] Si $E(X_n)=m$ et $V(X_n)\le C/n^2$, démontrer $X_n\to m$ en probabilité.

\section*{Exercice 23.5 -- Intervalle aléatoire}
[Raisonnement] Pour une fréquence $F_n$, justifier que l'intervalle aléatoire $[F_n-\varepsilon,F_n+\varepsilon]$ contient $p$ avec probabilité au moins $1-1/(4n\varepsilon^2)$.

\section*{Exercice 23.6 -- Attente géométrique}
[Approfondissement] Des essais indépendants de succès $p\in]0,1[$ sont répétés. Pour le rang $T$ du premier succès, calculer $P(T>n)$ et $P(T>m+n\mid T>m)$.

\chapter{Algorithmes : calculer, prouver et contrôler l’erreur}
\section*{Exercice 24.1 -- Invariant de boucle}
[Programmation] Écrire une boucle calculant $n!$ pour $n\ge0$ et justifier sa correction et sa terminaison.

\section*{Exercice 24.2 -- Coût de la dichotomie}
[Méthode] Sur un intervalle de longueur 3, combien de divisions suffisent pour une erreur du milieu au plus $10^{-6}$ ?

\section*{Exercice 24.3 -- Newton cyclique}
[Critique] Pour $f(x)=x^3-2x+2$, calculer deux itérations de Newton depuis 0. Quelle protection prévoir dans un programme ?

\section*{Exercice 24.4 -- Erreur des rectangles}
[Preuve] Pour $\ln2=\int_1^2dx/x$, donner un encadrement par $n$ rectangles et la largeur de cet encadrement.

\section*{Exercice 24.5 -- Monte-Carlo et erreur}
[Application] On estime $\pi$ par $\widehat\pi_n=4S_n/n$, où $S_n$ compte les points uniformes indépendants tombant dans le quart de disque unité. Majorer la probabilité d'une erreur au moins $\varepsilon$.

\section*{Exercice 24.6 -- Résidu et erreur}
[Raisonnement] Si $f(x)=10^{-9}(x-1000)$, un résidu $|f(0)|=10^{-6}$ signifie-t-il que 0 est proche de la racine ?

\chapter{Construire et calculer les nombres complexes}
\section*{Exercice 25.1 -- Calcul algébrique}
[Base] Simplifier $(2-3i)(1+2i)$ et $(1+i)/(2-i)$.

\section*{Exercice 25.2 -- Conjugué}
[Méthode] Résoudre $z+2\overline z=3-i$.

\section*{Exercice 25.3 -- Inverse et module}
[Preuve] Si $z\ne0$, démontrer $z^{-1}=\overline z/|z|^2$ et $|z^{-1}|=1/|z|$.

\section*{Exercice 25.4 -- Binôme complexe}
[Méthode] Calculer $(1+i)^4$ par deux méthodes.

\section*{Exercice 25.5 -- Absence d’ordre compatible}
[Raisonnement] Pourquoi ne peut-on munir $\mathbb C$ d'un ordre total compatible avec les opérations comme celui de $\R$ ?

\section*{Exercice 25.6 -- Inégalité triangulaire}
[Preuve] Montrer $|z+w|\le|z|+|w|$ et donner un cas où l'inégalité est stricte.

\chapter{Arguments, exponentielle imaginaire et trigonométrie}
\section*{Exercice 26.1 -- Forme trigonométrique}
[Base] Écrire $-1+i\sqrt3$ sous forme trigonométrique et exponentielle.

\section*{Exercice 26.2 -- Puissance}
[Méthode] Calculer $(\sqrt3+i)^6$.

\section*{Exercice 26.3 -- Addition trigonométrique}
[Preuve] Déduire $\cos(a+b)$ et $\sin(a+b)$ de $e^{i(a+b)}=e^{i a}e^{i b}$.

\section*{Exercice 26.4 -- Linéarisation}
[Méthode] Exprimer $\cos^3x$ comme combinaison de cosinus d'angles multiples.

\section*{Exercice 26.5 -- Somme de phases}
[Approfondissement] Pour $e^{i t}\ne1$, calculer $\sum_{k=0}^{n-1}e^{i kt}$. Que se passe-t-il si $t$ est multiple de $2\pi$ ?

\section*{Exercice 26.6 -- Argument d’un produit}
[Critique] Les arguments principaux de deux nombres sont $3\pi/4$ et $3\pi/4$. L'argument principal du produit est-il $3\pi/2$ dans $]-\pi,\pi]$ ?

\chapter{Polynômes, racines et factorisation dans le plan complexe}
\section*{Exercice 27.1 -- Second degré réel}
[Base] Résoudre $z^2-2z+5=0$ dans $\mathbb C$.

\section*{Exercice 27.2 -- Factorisation}
[Méthode] Factoriser $P(z)=z^3-1$ puis résoudre $P(z)=0$.

\section*{Exercice 27.3 -- Racine double}
[Raisonnement] Factoriser $z^3-3z+2$ sachant que 1 est racine. Combien y a-t-il de racines distinctes ?

\section*{Exercice 27.4 -- Racines conjuguées}
[Preuve] Montrer que si un polynôme à coefficients réels a pour racine $a+i b$, alors $a-i b$ est aussi racine.

\section*{Exercice 27.5 -- Nombre de racines}
[Preuve] Pourquoi un polynôme de degré au plus 4 s'annulant en cinq points distincts est-il identiquement nul ?

\section*{Exercice 27.6 -- Racine complexe d’un complexe}
[Approfondissement] Résoudre $z^2=3+4i$ en posant $z=x+i y$.

\chapter{Géométrie complexe et racines de l’unité}
\section*{Exercice 28.1 -- Distances}
[Base] Décrire les ensembles $|z-1|=2$ et $|z-1|=|z+1|$.

\section*{Exercice 28.2 -- Rotation}
[Méthode] Décrire la transformation $z'=i z+1-i$.

\section*{Exercice 28.3 -- Angle droit}
[Méthode] Pour $a=0$, $b=2$, $c=1+i$, vérifier que le triangle $ACB$ est rectangle en $C$.

\section*{Exercice 28.4 -- Racines quatrièmes}
[Base] Résoudre $z^4=1$ et calculer la somme des solutions.

\section*{Exercice 28.5 -- Racines d’un nombre non nul}
[Approfondissement] Résoudre $z^3=-8i$.

\section*{Exercice 28.6 -- Somme des racines de l’unité}
[Preuve] Pour $n\ge2$, démontrer que la somme des $n$ racines $n$-ièmes de l'unité est nulle. Interpréter géométriquement.

\chapter{Divisibilité, congruences et algorithme d’Euclide}
\section*{Exercice 29.1 -- Division euclidienne}
[Base] Effectuer la division euclidienne de $-17$ par 5.

\section*{Exercice 29.2 -- Dernier chiffre}
[Méthode] Déterminer le dernier chiffre de $7^{2026}$.

\section*{Exercice 29.3 -- Euclide}
[Base] Calculer $\operatorname{pgcd}(252,198)$.

\section*{Exercice 29.4 -- Divisibilité par 9}
[Preuve] Démontrer le critère de divisibilité par 9 d'un entier écrit en base dix.

\section*{Exercice 29.5 -- Simplification interdite}
[Critique] De $2x\equiv2y\pmod6$, peut-on conclure $x\equiv y\pmod6$ ?

\section*{Exercice 29.6 -- Carrés modulo 4}
[Raisonnement] Montrer qu'un carré entier est congru à 0 ou 1 modulo 4. En déduire que $x^2+y^2=4k+3$ n'a pas de solution entière.

\chapter{Bézout, Gauss et équations diophantiennes}
\section*{Exercice 30.1 -- Bézout explicite}
[Méthode] Trouver des entiers $u,v$ tels que $35u+22v=1$.

\section*{Exercice 30.2 -- Inverse modulaire}
[Base] Résoudre $7x\equiv3\pmod{11}$.

\section*{Exercice 30.3 -- Équation diophantienne}
[Méthode] Résoudre $6x+9y=15$ dans $\Z^2$.

\section*{Exercice 30.4 -- Aucune solution}
[Critique] L'équation $14x+21y=5$ admet-elle une solution entière ?

\section*{Exercice 30.5 -- Gauss et puissances}
[Preuve] Si $a$ et $b$ sont premiers entre eux, montrer que $a$ et $b^2$ le sont aussi.

\section*{Exercice 30.6 -- Deux congruences}
[Approfondissement] Résoudre $x\equiv2\pmod3$, $x\equiv3\pmod5$.

\chapter{Nombres premiers, Fermat et chiffrement}
\section*{Exercice 31.1 -- Tester un nombre premier}
[Base] Montrer que 97 est premier.

\section*{Exercice 31.2 -- Factorisation et pgcd}
[Méthode] Décomposer 360 et 252 en facteurs premiers ; en déduire leur pgcd et leur ppcm.

\section*{Exercice 31.3 -- Fermat}
[Méthode] Calculer le reste de $3^{100}$ modulo 7.

\section*{Exercice 31.4 -- Infinité des nombres premiers}
[Preuve] Si $p_1,\ldots,p_r$ est une liste finie de nombres premiers, montrer qu'il existe un nombre premier hors de cette liste.

\section*{Exercice 31.5 -- Chiffrement jouet}
[Application] Pour RSA avec $p=3,q=11,n=33,e=3,d=7$, chiffrer le message $m=2$ et le déchiffrer. Vérifier la relation entre $e,d$ et $(p-1)(q-1)$.

\section*{Exercice 31.6 -- Réciproque de Fermat}
[Critique] Le nombre $341=11\times31$ vérifie $2^{340}\equiv1\pmod{341}$. Le démontrer sans développer la puissance, puis conclure sur un test de primalité fondé sur une seule base.

\chapter{Matrices, applications linéaires et systèmes}
\section*{Exercice 32.1 -- Produit non commutatif}
[Base] Pour $A=\begin{pmatrix}1&1\\0&1\end{pmatrix}$ et $B=\begin{pmatrix}1&0\\1&1\end{pmatrix}$, calculer $AB$ et $BA$.

\section*{Exercice 32.2 -- Inverse et système}
[Méthode] Résoudre $2x+y=5$, $x+y=3$ à l'aide d'une matrice inverse.

\section*{Exercice 32.3 -- Produit nul}
[Critique] Donner deux matrices non nulles dont le produit est nul. Peut-on simplifier $AB=AC$ par une matrice $A$ quelconque ?

\section*{Exercice 32.4 -- Rotation matricielle}
[Preuve] Montrer que $R_t=\begin{pmatrix}\cos t&-\sin t\\\sin t&\cos t\end{pmatrix}$ conserve la norme d'un vecteur.

\section*{Exercice 32.5 -- Trois inconnues}
[Méthode] Résoudre $x+y+z=6$, $x-y+z=2$, $2x+y-z=1$.

\section*{Exercice 32.6 -- Application linéaire}
[Approfondissement] L'application $T(x,y)=(x+2y,3x-y)$ est-elle linéaire ? Écrire sa matrice. Comparer avec $S(x,y)=(x+1,y)$.

\chapter{Puissances de matrices et suites couplées}
\section*{Exercice 33.1 -- Nilpotence}
[Preuve] Si $N^2=0$, montrer que $(I+N)^n=I+nN$ pour tout entier $n\ge0$.

\section*{Exercice 33.2 -- Puissance par diagonalisation}
[Méthode] Pour $A=\begin{pmatrix}2&1\\1&2\end{pmatrix}$, utiliser les vecteurs $(1,1)$ et $(1,-1)$ pour calculer $A^n$.

\section*{Exercice 33.3 -- Suites couplées}
[Application] Résoudre $u_{n+1}=2u_n+v_n$, $v_{n+1}=u_n+2v_n$, avec $u_0=1,v_0=0$.

\section*{Exercice 33.4 -- Récurrence affine vectorielle}
[Méthode] Si $U_{n+1}=\frac12U_n+B$ avec $B=\binom12$ et $U_0=\binom00$, déterminer $U_n$ et sa limite.

\section*{Exercice 33.5 -- Fibonacci}
[Approfondissement] Pour $F_0=0,F_1=1,F_{n+2}=F_{n+1}+F_n$, écrire une récurrence matricielle et calculer $F_6$.

\section*{Exercice 33.6 -- Binôme matriciel}
[Critique] Pourquoi $(A+B)^2=A^2+2AB+B^2$ est-il faux en général ? Donner la formule correcte.

\chapter{Graphes, chemins et matrices d’adjacence}
\section*{Exercice 34.1 -- Degrés}
[Base] Un graphe simple non orienté a six sommets de degrés $1,2,2,3,3,3$. Combien possède-t-il d'arêtes ?

\section*{Exercice 34.2 -- Adjacence}
[Méthode] Écrire la matrice du chemin non orienté $A-B-C$, dans cet ordre. Calculer le nombre de marches de longueur 2 de $A$ à $A$ et de $A$ à $C$.

\section*{Exercice 34.3 -- Connexité}
[Raisonnement] Un graphe simple possède cinq sommets et deux composantes connexes de tailles 2 et 3. Quel est le nombre maximal d'arêtes ?

\section*{Exercice 34.4 -- Sommets de degré impair}
[Preuve] Montrer qu'un graphe fini non orienté possède un nombre pair de sommets de degré impair.

\section*{Exercice 34.5 -- Chemins et puissances}
[Preuve] Expliquer pourquoi $(M^{n+1})_{ij}=\sum_k(M^n)_{ik}M_{kj}$ compte les marches de longueur $n+1$.

\section*{Exercice 34.6 -- Orientation}
[Application] Dans le cycle orienté $A\to B\to C\to A$, combien de marches de longueur 2026 vont de $A$ à $B$ ?

\chapter{Chaînes de Markov et distributions invariantes}
\section*{Exercice 35.1 -- Convention de transition}
[Base] Pour $P=\begin{pmatrix}0{,}8&0{,}2\\0{,}3&0{,}7\end{pmatrix}$ et la distribution ligne $\mu_0=(1,0)$, calculer $\mu_1,\mu_2$.

\section*{Exercice 35.2 -- Distribution invariante}
[Méthode] Pour la matrice précédente, déterminer la distribution invariante et la limite de $\mu_n$.

\section*{Exercice 35.3 -- Périodicité}
[Critique] Pour $P=\begin{pmatrix}0&1\\1&0\end{pmatrix}$, trouver une distribution invariante et expliquer pourquoi la distribution issue de $(1,0)$ ne converge pas.

\section*{Exercice 35.4 -- Trois états}
[Méthode] Pour $P=\begin{pmatrix}1/2&1/2&0\\0&1/2&1/2\\1/2&0&1/2\end{pmatrix}$, vérifier que $(1/3,1/3,1/3)$ est invariante et calculer la distribution après deux pas depuis l'état 1.

\section*{Exercice 35.5 -- Probabilité d’un trajet}
[Application] Pour la matrice de l'exercice 35.1, partir de l'état 1 et calculer la probabilité du trajet $1,2,2,1$.

\section*{Exercice 35.6 -- État absorbant}
[Approfondissement] Pour $P=\begin{pmatrix}q&1-q\\0&1\end{pmatrix}$, $0<q<1$, partir du premier état. Calculer la probabilité d'y rester après $n$ pas et le temps moyen avant absorption.

\chapter{Fondements pour la licence : complétude, intégration et abstraction}
\section*{Exercice 36.1 -- Supremum et maximum}
[Base] Déterminer la borne supérieure de $A=\{1-1/n:n\ge1\}$. Est-ce un maximum ?

\section*{Exercice 36.2 -- Une hypothèse indispensable}
[Critique] Une fonction continue sur un intervalle est-elle toujours bornée et atteint-elle toujours ses bornes ?

\section*{Exercice 36.3 -- Convergence uniforme}
[Approfondissement] Pour $f_n(x)=x/n$ sur $[0,1]$, montrer la convergence uniforme vers zéro. Que devient cette affirmation sur $\R$ ?

\section*{Exercice 36.4 -- Noyau d’une dérivation}
[Raisonnement] Sur les polynômes de degré au plus 2, écrire la matrice de la dérivation dans la base $(1,X,X^2)$ et déterminer son noyau.

\section*{Exercice 36.5 -- Échange limite–intégrale}
[Approfondissement] Si $f_n,f$ sont continues sur $[a,b]$ et $f_n\to f$ uniformément, montrer $\int_a^bf_n\to\int_a^bf$.

\section*{Exercice 36.6 -- Audit d’une preuve}
[Synthèse courte] Un élève écrit : « $u_{n+1}=\sqrt{2+u_n}$, donc la limite vérifie $L=\sqrt{2+L}$ et vaut 2. » Identifier les étapes manquantes et compléter pour $u_0=0$.

\chapter{Problèmes de synthèse}
Les synthèses S1 à S4 et S8 mobilisent la spécialité et ses approfondissements signalés ; S5 à S7 mobilisent aussi les mathématiques expertes.
\section{S1 -- Héron : convergence et efficacité}
Pour $a>0$, $u_0>0$, poser $u_{n+1}=\tfrac12(u_n+a/u_n)$.
\begin{enumerate}
\item Montrer que la suite est définie et que $u_n\ge\sqrt a$ pour $n\ge1$.
\item Établir la décroissance à partir du rang 1, puis la convergence vers $\sqrt a$.
\item Démontrer $u_{n+1}-\sqrt a=(u_n-\sqrt a)^2/(2u_n)$.
\item Pour $a=2,u_0=2$, donner un encadrement rationnel de $\sqrt2$ après trois itérations. Expliquer comment certifier une précision sans connaître la racine.
\item Relier la méthode à Newton et expliquer pourquoi l'exemple ne justifie pas la convergence de Newton pour toute fonction.
\end{enumerate}

\section{S2 -- Géométrie d’un tétraèdre}
Dans un repère orthonormé, $O=(0,0,0)$, $A=(2,0,0)$, $B=(0,3,0)$, $C=(0,0,6)$ ; les coordonnées sont en centimètres.
\begin{enumerate}
\item Déterminer une équation du plan $(ABC)$ et un vecteur normal.
\item Calculer le projeté orthogonal $H$ de $O$ sur ce plan, puis $OH$.
\item Calculer le volume de $OABC$ en prenant $OAB$ pour base, puis l'aire de $ABC$ en utilisant la hauteur $OH$.
\item Montrer que la sphère de diamètre $[OC]$ a équation $x^2+y^2+z^2-6z=0$.
\item Déterminer son intersection avec le plan $z=3$ ; expliquer pourquoi aucune figure seule ne prouve ces résultats.
\end{enumerate}

\section{S3 -- Logarithme, intégrale et sommes}
Pour $n\ge1$, poser $L_n=\frac1n\sum_{k=0}^{n-1}(1+k/n)^{-1}$ et $R_n=\frac1n\sum_{k=1}^{n}(1+k/n)^{-1}$.
\begin{enumerate}
\item Justifier $R_n\le\ln2\le L_n$ et calculer $L_n-R_n$.
\item En déduire un algorithme certifiant une erreur au plus $10^{-4}$ par le milieu de l'encadrement.
\item Poser $J_n=\int_0^1x^n/(1+x)\,dx$. Montrer $0\le J_n\le1/(n+1)$ et $J_n+J_{n+1}=1/(n+1)$.
\item En déduire une expression de $\ln2$ par une somme alternée avec une borne de reste.
\item Comparer approximation, encadrement prouvé et valeur affichée en décimal.
\end{enumerate}

\section{S4 -- Surbooking : modèle et décision}
Une salle possède 100 places. Pour une simulation pédagogique, chaque personne inscrite vient indépendamment avec probabilité $0{,}9$ ; on inscrit $n\ge100$ personnes.
\begin{enumerate}
\item Définir la loi du nombre $X_n$ de présents, sa moyenne et sa variance.
\item Écrire exactement $P(X_n>100)$ et proposer un algorithme cherchant le plus grand $n$ tel que ce risque soit au plus $0{,}05$.
\item Justifier que ce risque croît avec $n$ et tend vers $1$, afin de justifier l'arrêt de la recherche.
\item Pour $n=105$, donner une majoration par Tchebychev et dire si elle décide la contrainte de $5\,\%$.
\item Expliquer quelles hypothèses devraient être examinées avant une application réelle.
\end{enumerate}

\section{S5 -- Complexes, rotations et polygones}
Poser $\omega=e^{2i\pi/5}$ et $A_k$ d'affixe $\omega^k$, $0\le k\le4$.
\begin{enumerate}
\item Prouver que les cinq points sont distincts et forment un pentagone régulier.
\item Calculer leur somme et leur barycentre.
\item Montrer $|1-e^{it}|=2|\sin(t/2)|$ ; en déduire le côté et la diagonale du pentagone.
\item Poser $s=\omega+\omega^{-1}$. À partir de la somme des racines, montrer $s^2+s-1=0$ et déterminer $\cos(2\pi/5)$.
\item Écrire la matrice de la rotation envoyant $A_k$ sur $A_{k+1}$ et montrer que sa cinquième puissance est l'identité.
\end{enumerate}

\section{S6 -- RSA : preuve et limites}
Prendre deux nombres premiers distincts $p=5,q=11$, $N=pq$, $e=3$.
\begin{enumerate}
\item Trouver $d\in\{1,\ldots,39\}$ tel que $ed\equiv1\pmod{40}$.
\item Chiffrer $m=7$ par $c\equiv m^e\pmod N$, puis déchiffrer par $c^d\pmod N$ sans développer de très grands entiers.
\item Prouver que $m^{ed}\equiv m$ modulo $5$ puis modulo $11$ pour tout entier $m$, y compris s'il est divisible par l'un de ces nombres.
\item En déduire la correction modulo $55$ à l'aide de Gauss.
\item Pourquoi cet exemple ne fournit-il pas un système sécurisé prêt à l'emploi ?
\end{enumerate}

\section{S7 -- Trois états et convergence}
Une chaîne de Markov possède la matrice de transition
$P=\begin{pmatrix}1/2&1/4&1/4\\1/4&1/2&1/4\\1/4&1/4&1/2\end{pmatrix}$, avec des distributions lignes.
\begin{enumerate}
\item Vérifier que $P$ est une matrice de transition et que $\pi=(1/3,1/3,1/3)$ est invariante.
\item Noter $J$ la matrice dont tous les coefficients valent $1$ et $Q=J/3$. Montrer $Q^2=Q$ et $P=Q+\tfrac14(I-Q)$.
\item En déduire $P^n=Q+4^{-n}(I-Q)$ pour $n\ge0$.
\item Déterminer la distribution au rang $n$ pour une distribution initiale quelconque $\mu_0$ ; prouver convergence et unicité de l'invariante.
\item Trouver un rang garantissant que chaque coordonnée diffère de $1/3$ d'au plus $10^{-3}$, quel que soit $\mu_0$.
\end{enumerate}

\section{S8 -- Une fonction définie par une intégrale}
Pour $x\ge0$, poser $F(x)=\int_0^xe^{-t^2}\,dt$.
\begin{enumerate}
\item Justifier que $F$ est dérivable, strictement croissante, puis concave sur $[0,+\infty[$.
\item Pour $x\in[0,1]$, montrer $x-x^3/3\le F(x)\le x$.
\item Pour $x\ge1$, utiliser $t^2\ge t$ sur $[1,x]$ pour majorer $F(x)$ indépendamment de $x$. En déduire l'existence d'une limite finie $L$.
\item Montrer $0\le L-F(x)\le e^{-x}$ pour $x\ge1$.
\item Expliquer pourquoi cela ne donne pas une primitive élémentaire de $e^{-x^2}$, ni la valeur exacte de $L$. Relier la question à une future théorie des densités.
\end{enumerate}

\end{document}
